\subsection{A concrete failure of the naive 245-state compilation}
The source association certificate does not by itself justify changing only
the two rows of old state 178 and deleting old state 208.  The following exact
finite computation identifies a failure of that proposed source-to-binary
instruction contract.  It does not establish inequivalence of blank-tape
halting, and it does not exclude a different simulation invariant.

Use original state labels, dispatcher 96, head coordinate zero, marker at
$-1$, and header at 40.  The program counter is
\[
 p=\mathtt{11101001001010100111000000}.
\]
The primary physical registers are $(0,1,0,3,5,8,13,21,34)$ and the spare list
is $[0]$.  Encode each register by $B(v)=0\,1^{v+1}$, beginning at coordinate
41, and set all other cells to zero.  The declared source instruction is
decrement of logical register 1, namely physical register 2, whose value is
zero.  Its prescribed successor is
\[
 p_{\mathrm{expected}}=\mathtt{1110100100101010011100001}.
\]

After 23 transitions the machine is in state 218 at coordinate 23 reading
zero.  It moves left to state 154 at coordinate 22, also reading zero, and
then returns to coordinate 23 in state 178.  At transition 26 the modified
row sends the computation to state 16 instead of state 23.  This is a
backward use of the decoder: the preceding bit, already read, is used again
to select a continuation.  The direct source-tree rotation does not describe
this additional use.

The first selected register is nevertheless the intended physical register
2: state 31 reaches the register boundary at transition 149, and state 95
starts its decrement operation there after 267 transitions.  After 378
transitions the machine returns to dispatcher 96 with head zero, the same
header, and the complete unchanged register and spare file, but with
\[
 p_{\mathrm{actual}}=\mathtt{11101001001010100111000001}.
\]
The full expected and actual tapes differ at exactly two coordinates: cell
24 should contain one but contains zero, and cell 25 should contain zero but
contains one.  Thus the error is in the program-counter update, not in the
first register selection.  At this returned counter, the declared rotated
source decodes increment of logical register 6 (physical register 8, selector
state 69), whereas the literal machine next enters selector state 60 for
physical register 0 after 531 transitions.

The accompanying independent dense-tape checker records these states and
the complete-frame discrepancy in \texttt{rotation\_counterexample.json}.
Failure of this local correspondence is a reason to reject this particular
compilation certificate; it is not a claim that the two blank-tape machines
have different halting behavior.
